\documentclass[12pt,a4paper]{article}

\usepackage{geometry}
\geometry{top=2.2cm,bottom=2.3cm,left=2.0cm,right=2.0cm}
\usepackage[most]{tcolorbox}
\usepackage{xcolor}
\usepackage{tikz}
\usetikzlibrary{positioning,arrows.meta,calc}
\usepackage{fontawesome5}
\usepackage{setspace}
\usepackage{enumitem}
\usepackage{hyperref}
\usepackage{graphicx}
\usepackage{float}
\usepackage{listings}
\usepackage{fancyhdr}
\usepackage{titlesec}
\usepackage{array}
\usepackage{tabularx}
\usepackage{caption}
\usepackage{booktabs}
\usepackage{colortbl}
\usepackage{xepersian}

\IfFontExistsTF{B Nazanin}{\settextfont{B Nazanin}}%
  {\IfFontExistsTF{Vazirmatn}{\settextfont{Vazirmatn}}{\settextfont{Amiri}}}
\IfFontExistsTF{Liberation Serif}{\setlatintextfont{Liberation Serif}}%
  {\setlatintextfont{Times New Roman}}
\IfFontExistsTF{JetBrainsMono Nerd Font}{\setmonofont{JetBrainsMono Nerd Font}}%
  {\setmonofont{DejaVu Sans Mono}}
\setstretch{1.32}

\definecolor{ink}{HTML}{20242A}
\definecolor{paper}{HTML}{FBFCFE}
\definecolor{midnight}{HTML}{263159}
\definecolor{navy}{HTML}{1E3A5F}
\definecolor{teal}{HTML}{0F766E}
\definecolor{green}{HTML}{10B981}
\definecolor{blue}{HTML}{3B82F6}
\definecolor{sky}{HTML}{EAF3FF}
\definecolor{mint}{HTML}{ECFDF3}
\definecolor{amber}{HTML}{F59E0B}
\definecolor{gold}{HTML}{E9B949}
\definecolor{softgold}{HTML}{FFF7D6}
\definecolor{rose}{HTML}{F43F5E}
\definecolor{lavender}{HTML}{F3E8FF}
\definecolor{violet}{HTML}{7C3AED}
\definecolor{softgray}{HTML}{F5F7FB}
\definecolor{linegray}{HTML}{D8DEE9}

\hypersetup{colorlinks=true,linkcolor=navy,urlcolor=teal,citecolor=teal}
\pagecolor{paper}
\color{ink}
\setlength{\parindent}{0pt}
\setlength{\parskip}{0.5em}
\setlength{\emergencystretch}{3em}

\setlist[itemize]{
  label=\textcolor{teal}{\scriptsize\faCircle},
  itemsep=0.22em,topsep=0.3em,rightmargin=1.35em
}
\setlist[enumerate]{itemsep=0.24em,topsep=0.3em,rightmargin=1.35em}

\lstdefinestyle{codestyle}{
  basicstyle=\ttfamily\footnotesize,
  breaklines=true,
  frame=single,
  backgroundcolor=\color{softgray},
  rulecolor=\color{blue},
  numbers=none,
  xleftmargin=0pt,xrightmargin=0pt,
  aboveskip=0.5em,belowskip=0.5em
}
\lstset{style=codestyle}

\newtcolorbox{tocbox}{
  enhanced,breakable,colback=white,colframe=teal,boxrule=0.9pt,arc=5pt,
  right=12pt,left=12pt,top=10pt,bottom=10pt,
  title={\textcolor{white}{\faListUl\hspace{0.45em} فهرست مطالب}},
  colbacktitle=teal,coltitle=white,fonttitle=\bfseries,
  attach boxed title to top right={yshift=-2mm,xshift=-2mm},
  boxed title style={arc=4pt,boxrule=0pt}
}

\newtcolorbox{infobox}[1]{
  enhanced,breakable,colback=sky,colframe=blue,boxrule=0.8pt,arc=5pt,
  right=10pt,left=10pt,top=8pt,bottom=8pt,
  title={\textcolor{white}{\faInfoCircle\hspace{0.45em} #1}},
  colbacktitle=blue,coltitle=white,fonttitle=\bfseries
}

\newtcolorbox{resultbox}[1]{
  enhanced,breakable,colback=mint,colframe=green,boxrule=0.8pt,arc=5pt,
  right=10pt,left=10pt,top=8pt,bottom=8pt,
  title={\textcolor{white}{\faCheckCircle\hspace{0.45em} #1}},
  colbacktitle=green,coltitle=white,fonttitle=\bfseries
}

\newtcolorbox{warningbox}[1]{
  enhanced,breakable,colback=softgold,colframe=amber,boxrule=0.8pt,arc=5pt,
  right=10pt,left=10pt,top=8pt,bottom=8pt,
  title={\textcolor{white}{\faExclamationTriangle\hspace{0.45em} #1}},
  colbacktitle=amber,coltitle=white,fonttitle=\bfseries
}

\newtcolorbox{scorebox}[1]{
  enhanced,breakable,colback=lavender,colframe=violet,boxrule=0.8pt,arc=5pt,
  right=10pt,left=10pt,top=8pt,bottom=8pt,
  title={\textcolor{white}{\faStar\hspace{0.45em} #1}},
  colbacktitle=violet,coltitle=white,fonttitle=\bfseries
}

\titleformat{\section}
  {\normalfont\Large\bfseries\color{midnight}}{}{0pt}
  {\textcolor{gold}{\faBookOpen}\hspace{0.55em}}
  [\vspace{0.2em}{\color{gold}\titlerule[0.9pt]}]
\titleformat{\subsection}
  {\normalfont\large\bfseries\color{teal}}{}{0pt}
  {\faBookmark\hspace{0.45em}}
\titleformat{\subsubsection}
  {\normalfont\normalsize\bfseries\color{navy}}{}{0pt}
  {\faAngleLeft\hspace{0.4em}}
\titlespacing*{\section}{0pt}{1.15em}{0.3em}
\titlespacing*{\subsection}{0pt}{0.85em}{0.18em}

\pagestyle{fancy}
\setlength{\headheight}{14.5pt}
\fancyhf{}
\fancyhead[R]{\small پروژه پایانی رایانش ابری ـ فاز دوم}
\fancyhead[L]{\small حسینیان ـ قویدل}
\fancyfoot[C]{\thepage}
\renewcommand{\headrulewidth}{0.4pt}

\newcommand{\safelogo}[2]{%
  \IfFileExists{#1}{\includegraphics[width=#2]{#1}}%
  {\fbox{\parbox[c][1.8cm][c]{2.0cm}{\centering لوگو}}}%
}
\newcommand{\codefile}[1]{\lr{\nolinkurl{#1}}}
\newcommand{\eng}[1]{\lr{#1}}

\newcommand{\makecover}{
\begin{titlepage}
\thispagestyle{empty}
\begin{tikzpicture}[remember picture,overlay]
  \fill[midnight] (current page.south west) rectangle (current page.north east);
  \fill[teal] (current page.north west) -- ++(0,-8.3)
    .. controls +(4.8,-1.1) and +(-4.8,1.3) ..
    ([yshift=-5.5cm]current page.north east)
    -- (current page.north east) -- cycle;
  \fill[blue] ([xshift=1.0cm,yshift=1.0cm]current page.south west) circle (3.3cm);
  \fill[gold] ([xshift=-2.25cm,yshift=-2.2cm]current page.north east) circle (3.25cm);
  \fill[softgold] ([xshift=-2.25cm,yshift=-2.2cm]current page.north east) circle (2.25cm);
  \draw[gold,line width=1.1pt]
    ([xshift=1.55cm,yshift=-10.7cm]current page.north west) --
    ([xshift=-1.55cm,yshift=-10.7cm]current page.north east);
\end{tikzpicture}

\vspace*{1.1cm}
\begin{center}
\begin{minipage}{0.18\textwidth}
\centering\safelogo{ce.png}{2.05cm}
\end{minipage}
\begin{minipage}{0.56\textwidth}
\centering
{\Large\textcolor{softgold}{دانشگاه صنعتی امیرکبیر}}\\[0.15cm]
{\large\textcolor{white}{دانشکده مهندسی کامپیوتر}}
\end{minipage}
\begin{minipage}{0.18\textwidth}
\centering\safelogo{aut.png}{2.15cm}
\end{minipage}\\[0.8cm]

{\Huge\bfseries\textcolor{gold}{فاز دوم پروژه پایانی}}\\[0.4cm]
{\LARGE\textcolor{white}{استقرار Kubernetes و خودکارسازی GitLab CI/CD}}\\[0.25cm]
{\Large\textcolor{softgold}{تحلیل لاگ جام جهانی}}\\[0.65cm]
{\textcolor{white}{\rule{0.72\textwidth}{0.8pt}}}\\[0.75cm]

\begin{tcolorbox}[
  enhanced,width=0.80\textwidth,colback=white,colframe=gold,
  boxrule=1pt,arc=5pt,right=18pt,left=18pt,top=13pt,bottom=13pt
]
\centering
{\Large\bfseries\textcolor{midnight}{مبانی رایانش ابری}}\\[0.3cm]
{\large\textcolor{teal}{نیمسال دوم ۱۴۰۴--۱۴۰۵}}\\[0.3cm]
{\large\bfseries\textcolor{midnight}{گزارش طراحی، پیاده‌سازی، خطایابی و آزمون}}
\end{tcolorbox}

\vspace{0.5cm}
\begin{tcolorbox}[
  enhanced,width=0.80\textwidth,colback=softgray,colframe=teal,
  boxrule=0.9pt,arc=5pt,right=18pt,left=18pt,top=13pt,bottom=13pt
]
\centering
{\large\bfseries\textcolor{midnight}{مشخصات گروه}}\\[0.25cm]
{\large\bfseries سید حسام الدین حسینیان}\quad{\large ۴۰۲۳۱۹۰۱}\\[0.18cm]
{\large\bfseries فروزان قویدل}\quad{\large ۴۰۲۳۱۰۶۴}\\[0.35cm]
{\large استاد درس: دکتر جوادی}\\
{\large تدریسیاران: مهرداد شیخ‌عباسی و محمدهادی ستاک}
\end{tcolorbox}

\vfill
{\large\textcolor{softgold}{تابستان ۱۴۰۵}}
\end{center}
\end{titlepage}
}

\begin{document}
\makecover
\setcounter{page}{1}

\begin{tocbox}
\small
\setcounter{tocdepth}{2}
\tableofcontents
\end{tocbox}

\clearpage
\section{خلاصه اجرایی}

در فاز دوم، مسیر دستی فاز اول به یک سامانه واقعی روی Kubernetes منتقل شد. یک کلاستر
\eng{kind} با namespace ثابت \codefile{cc-project} ساخته شد. GitLab CE، GitLab
Runner، Registry خصوصی، سه backend، Nginx، HDFS، Jobهای ترافیک، MapReduce،
اعتبارسنجی و تولید گزارش همگی داخل همان کلاستر اجرا شدند. کد پروژه در GitLab
محلی push شد و Runner با \eng{Kubernetes executor} همه مراحل را داخل podهای
موقت اجرا کرد.

در اجرای نهایی، ۱۰۰٬۰۰۰ درخواست فقط از Nginx عبور کرد. فایل access log شامل
دقیقاً ۱۰۰٬۰۰۰ شناسه یکتا بود و مجموع سه فایل backend نیز دقیقاً ۱۰۰٬۰۰۰ رکورد
شد. پنج Job واقعی Hadoop Streaming روی HDFS اجرا شدند؛ ۱۳ خروجی الزامی با
validator کم‌حافظه تأیید شد؛ داشبورد HTML ساخته و با NodePort روی پورت ۹۰۰۰
منتشر شد.

\begin{resultbox}{وضعیت نهایی بخش اجباری}
کلاستر، GitLab، Runner، Registry، پایپ‌لاین هشت‌مرحله‌ای، ترافیک ۱۰۰هزار
درخواست، پنج Job Hadoop، اعتبارسنجی خروجی و داشبورد قابل مشاهده همگی به‌صورت
واقعی اجرا شده‌اند. اجرای نهایی پایپ‌لاین با شناسه ۱۶ و commit
\codefile{f507f6c3} در ۳۷۲ ثانیه با وضعیت موفق تمام شد. هیچ Secret یا token
در مخزن و بسته تحویلی قرار نگرفته است.
\end{resultbox}

\section{خواندن دستورکار و تبدیل آن به برنامه اجرا}

فایل ۲۵ صفحه‌ای فاز دوم ابتدا کامل خوانده شد. سپس الزام‌ها به ترتیب وابستگی
مرتب شدند. نتیجه این بود که اجرای درست باید از زیرساخت شروع شود، بعد GitLab و
Runner بالا بیایند، سپس build و deploy انجام شود و در پایان ترافیک، تحلیل و
انتشار اجرا شوند. ترتیب کاری استفاده‌شده چنین بود:

\begin{enumerate}
  \item بررسی کد فاز اول، template گزارش و منابع میزبان؛
  \item ساخت کپی مستقل فاز دوم و نگه‌داشتن خروجی فاز اول بدون تغییر؛
  \item طراحی namespace، PVC، Serviceها، Deploymentها و Jobها؛
  \item ساخت کلاستر \eng{kind} و کنترل سلامت node؛
  \item راه‌اندازی Registry، GitLab CE و Runner داخل Kubernetes؛
  \item ایجاد پروژه GitLab، ثبت Runner و push مخزن؛
  \item طراحی پایپ‌لاین test تا publish و build بدون Docker daemon با Kaniko؛
  \item اجرای ترافیک واقعی و رفع اشکال مسیرهای PVC؛
  \item اجرای Hadoop Streaming، اعتبارسنجی ۱۳ خروجی و ساخت dashboard؛
  \item ثبت همه خطاهای واقعی، اندازه‌گیری منابع و آماده‌سازی بسته تحویل.
\end{enumerate}

\subsection{وضعیت منابع در شروع}

میزبان ۷٫۴ گیگابایت RAM و ۱۶ پردازنده منطقی داشت. مقدار حافظه available در
بخش‌های مختلف معمولاً بین ۱٫۰ تا ۱٫۸ گیگابایت بود. به همین علت از ابتدا سه
تصمیم گرفته شد: Runner فقط یک Job هم‌زمان اجرا کند، HDFS تنها در مرحله MapReduce
روشن باشد و buildهای سنگین به‌صورت ترتیبی انجام شوند.

\begin{table}[H]
\centering
\caption{بودجه منابع بخش‌های پرمصرف}
\small
\begin{tabularx}{\textwidth}{>{\bfseries}X p{3.0cm} p{3.0cm}}
\toprule
بخش & request & limit \\
\midrule
GitLab CE & ۵۰۰m / ۱۵۰۰MiB & ۲ CPU / ۳۲۰۰MiB \\
GitLab Runner & ۵۰m / ۶۴MiB & ۵۰۰m / ۳۸۴MiB \\
MapReduce & ۲۵۰m / ۳۸۴MiB & ۱ CPU / ۱۲۰۰MiB \\
NameNode & ۱۰۰m / ۲۵۶MiB & ۱ CPU / ۸۰۰MiB \\
DataNode & ۱۰۰m / ۲۲۴MiB & ۶۰۰m / ۶۰۰MiB \\
Spark اختیاری & ۲۵۰m / ۵۱۲MiB & ۲ CPU / ۱۶۰۰MiB \\
\bottomrule
\end{tabularx}
\end{table}

requestها برای زمان‌بندی محافظه‌کارانه کوچک نگه داشته شدند، اما limitها جلوی
مصرف کنترل‌نشده را می‌گیرند. HDFS و Spark هم‌زمان اجرا نمی‌شوند. هنگام Spark
نیز GitLab و مسیر سرویس‌دهی موقتاً scale-down می‌شوند؛ بنابراین جمع limitهای
جدول به معنی رزرو هم‌زمان حافظه نیست.

\clearpage
\section{معماری نهایی سامانه}

\begin{figure}[H]
\centering
\IfFileExists{architecture.png}
  {\includegraphics[width=\textwidth,height=0.82\textheight,keepaspectratio]{architecture.png}}
  {\includegraphics[width=\textwidth]{../code_phase2/code/docs/architecture.png}}
\caption{گراف معماری نهایی فاز دوم}
\end{figure}

مسیر اصلی از چپ به راست در شکل دیده می‌شود. Developer کد را به GitLab داخلی
می‌فرستد. GitLab یک pipeline هشت‌مرحله‌ای ایجاد می‌کند و Runner با
\eng{Kubernetes executor} برای هر Job یک pod موقت می‌سازد. Kaniko imageها را
بدون دسترسی به Docker daemon در Registry خصوصی می‌گذارد. Deploymentها image
همان commit را pull می‌کنند. Job ترافیک فقط با Service داخلی Nginx صحبت می‌کند.
Nginx و backendها JSON log را در PVC مشترک می‌نویسند. در مرحله تحلیل،
NameNode و DataNode از صفر به یک scale می‌شوند، Job Hadoop خروجی را روی همان PVC
قرار می‌دهد و پس از پایان دوباره به صفر scale می‌شوند. validator و dashboard
آخرین مصرف‌کنندگان داده هستند.

\begin{table}[H]
\centering
\caption{اجزای اصلی معماری}
\small
\begin{tabularx}{\textwidth}{>{\bfseries}p{3.2cm}X p{3.2cm}}
\toprule
جزء & نقش & نوع Kubernetes \\
\midrule
GitLab CE & نگهداری مخزن و مدیریت CI/CD & Deployment + Service + PVC \\
GitLab Runner & اجرای Jobها با Kubernetes executor & Deployment + RBAC + Secret \\
Registry خصوصی & نگهداری imageهای هر commit و cache کانیکو & Deployment + ClusterIP + PVC \\
سه backend & پاسخ‌گویی و نوشتن service log & سه Deployment و سه Service \\
Nginx Gateway & تنها ورودی ترافیک و تولید access log & Deployment + NodePort \\
project-data & لاگ‌ها، خروجی Hadoop و dashboard & PVC دو گیگابایتی \\
HDFS & ذخیره ورودی و خروجی Streaming & دو StatefulSet و دو Service \\
MapReduce & اجرای پنج مرحله تحلیل & Job \\
Report Web & انتشار HTML و JSON & Deployment + NodePort \\
\bottomrule
\end{tabularx}
\end{table}

\section{پیاده‌سازی زیرساخت Kubernetes}

\subsection{ساخت کلاستر و namespace}

فایل \codefile{k8s/kind-config.yaml} پورت‌های NodePort لازم را از node به میزبان
منتقل می‌کند. namespace پروژه \codefile{cc-project} است. اسکریپت
\codefile{scripts/phase2/create\_cluster.sh} idempotent نوشته شد: اگر کلاستر
وجود داشته باشد آن را پاک نمی‌کند، context درست را انتخاب می‌کند، تا Ready شدن
node منتظر می‌ماند و سپس storage را apply می‌کند.

containerd در namespace شبکه node اجرا می‌شود و DNS داخلی podها را ندارد.
برای اینکه نام \codefile{registry.cc-project.svc} هنگام image pull نیز resolve
شود، اسکریپت ساخت کلاستر نگاشت ثابت \codefile{10.96.0.50} را داخل hosts نود
ثبت می‌کند. این مورد پس از مشاهده توقف image pull به اسکریپت اضافه شد.

\subsection{Storage و تفکیک داده}

چهار PVC مستقل تعریف شد:

\begin{itemize}
  \item \codefile{project-data} با ظرفیت ۲Gi برای لاگ و خروجی تحلیل؛
  \item \codefile{registry-data} با ظرفیت ۳Gi برای imageها و cache؛
  \item \codefile{gitlab-data} با ظرفیت ۸Gi برای repository و پایگاه داده؛
  \item \codefile{gitlab-config} با ظرفیت ۱Gi برای کلیدها و تنظیمات GitLab.
\end{itemize}

جداکردن \codefile{/etc/gitlab} از data یک اصلاح مهم بود. در نسخه اولیه config
روی \eng{emptyDir} قرار داشت. پس از جایگزینی pod، کلیدهای encryption دوباره
ساخته شدند و GitLab هنگام خواندن tokenهای قبلی خطای
\eng{OpenSSL::Cipher} داد. داده GitLab یک بار در محیط آزمایش بازسازی و پس از
آن config روی PVC پایدار قرار گرفت.

\subsection{ConfigMap، Secret و RBAC}

پیکربندی Nginx و Report Web در ConfigMap قرار گرفت. گذرواژه اولیه GitLab،
authentication token مربوط به Runner و API token کوتاه‌عمر در Secret یا فایل
محلی با mode برابر ۰۶۰۰ نگهداری شدند و مسیر
\codefile{.phase2-secrets} در gitignore قرار گرفت.

ServiceAccount مربوط به Runner فقط مجوزهای لازم برای pod، Job، Deployment،
StatefulSet، Service، ConfigMap و PVC در namespace پروژه را دارد. برای فرمان
\codefile{get namespace/cc-project} یک ClusterRole محدود با
\eng{resourceNames} تعریف شد. آزمون RBAC نشان داد get روی همان namespace مجاز
و list همه namespaceها غیرمجاز است.

\section{انتقال سرویس‌ها و Nginx}

\subsection{سه backend}

سرویس‌های match، team و stadium از کد فاز اول استفاده می‌کنند، اما اکنون هر
کدام Deployment و ClusterIP Service مستقل دارند. برای هر container request و
limit مشخص شد. probeهای readiness و liveness روی \codefile{/health} هستند.

فایل‌های لاگ باید در این مسیرهای مشترک نوشته شوند:

\begin{latin}
\begin{lstlisting}
/project-data/data/service_logs/match_service.log
/project-data/data/service_logs/team_service.log
/project-data/data/service_logs/stadium_service.log
\end{lstlisting}
\end{latin}

برنامه‌ها متغیر \codefile{SERVICE\_LOG\_PATH} را می‌خوانند. در یک اجرای واقعی
مشخص شد manifest اشتباهاً متغیر \codefile{LOG\_PATH} می‌فرستاد؛ به همین دلیل
حدود ۲۹MB لاگ داخل filesystem موقت سه container باقی مانده بود. نام متغیر و
مسیر هر سه Deployment اصلاح شد و یک regression check به تست‌های CI اضافه شد.

\subsection{Gateway و مسیر لاگ}

Nginx تنها درگاه عمومی API است. routeهای \codefile{/api/matches}،
\codefile{/api/teams} و \codefile{/api/stadiums} به Service داخلی متناظر
proxy می‌شوند. headerهای request ID، کشور و scenario بدون تغییر به backend
می‌رسند. health check از access log خارج شده است تا داده تحلیل فقط شامل
درخواست‌های پروژه باشد.

access log با فرمت JSON و buffer محدود در مسیر زیر ذخیره می‌شود:

\begin{latin}
\begin{lstlisting}
/project-data/data/nginx/nginx_access.log
\end{lstlisting}
\end{latin}

در اجرای اول فاز دو، Nginx فایل را در
\codefile{/project-data/nginx/nginx\_access.log} می‌نوشت ولی Hadoop مسیر
\codefile{data/nginx} را می‌خواند. همان اجرا پیش از MapReduce متوقف شد، فایل
اشتباه پاک شد، ConfigMap و initContainer اصلاح شدند و ترافیک از ابتدا تکرار شد.

\section{تولید ترافیک به شکل Kubernetes Job}

مولد ترافیک یک Job با ۳۲ worker محدود است. مقصد آن فقط
\codefile{http://nginx-gateway:80} است و هیچ آدرس backend در container وجود
ندارد. seed ثابت ۱۴۰۵ باعث می‌شود توزیع درخواست‌ها تکرارپذیر باشد. پنج scenario
عادی، popular، invalid، server error و slow تولید می‌شوند.

پیش از هر اجرا، stage ترافیک این ترتیب را انجام می‌دهد:

\begin{enumerate}
  \item replicaهای Nginx و سه backend را صفر می‌کند؛
  \item تا حذف pod قدیمی Nginx منتظر می‌ماند تا buffer قبلی باز نماند؛
  \item Job پاک‌سازی PVC را اجرا می‌کند و فایل‌ها را صفر می‌کند؛
  \item Deploymentها را دوباره یک می‌کند و rollout هر چهار مورد را کنترل می‌کند؛
  \item Job ترافیک را با image tag همان commit می‌سازد؛
  \item تا condition برابر Complete منتظر می‌ماند و log خلاصه را چاپ می‌کند.
\end{enumerate}

\begin{resultbox}{شمارش اجرای مبنا}
فایل Nginx دقیقاً ۱۰۰٬۰۰۰ خط و ۱۰۰٬۰۰۰ request ID یکتا داشت. فایل‌های match،
team و stadium به‌ترتیب ۲۸٬۱۶۲، ۴۹٬۶۱۹ و ۲۲٬۲۱۹ خط داشتند؛ جمع آن‌ها نیز
دقیقاً ۱۰۰٬۰۰۰ است. حجم access log حدود ۲۸٫۷MiB و مجموع service logها حدود
۲۷٫۴MiB بود.
\end{resultbox}

برای خطاهای گذرای شبکه، Job 1 رکوردها را با کلید ترکیبی source و request ID
مرتب و duplicate احتمالی را حذف می‌کند. یک تست جداگانه دو نسخه از یک درخواست
را به mapper/reducer می‌دهد و خروجی یک رکورد را کنترل می‌کند.

\section{GitLab، Runner و Registry داخل کلاستر}

\subsection{GitLab CE}

GitLab CE نسخه \eng{17.11.7-ce.0} داخل Deployment اجرا شد. قابلیت‌های غیرضروری
مانند monitoring داخلی برای کاهش مصرف خاموش شدند. readiness probe از curl
داخلی GitLab روی \codefile{127.0.0.1/-/readiness} استفاده می‌کند. دسترسی میزبان
از \codefile{http://localhost:8929} فراهم است، ولی Runner و سرویس‌های داخل
کلاستر از DNS داخلی \codefile{gitlab.cc-project.svc} استفاده می‌کنند.

اولین pull image GitLab حدود ۶ دقیقه و ۸ ثانیه طول کشید و image فشرده حدود
۱٫۷۸GB بود. هنگام migration، مصرف node حدود ۲٫۰GiB و CPU به‌طور لحظه‌ای نزدیک
۱۰۰ درصد یک هسته شد. پس از پایدار شدن، GitLab بخش اصلی حافظه را به Puma،
Sidekiq و PostgreSQL اختصاص داد.

\subsection{GitLab Runner}

Runner نسخه \eng{17.11.1} در همان namespace ثبت شد. executor برابر
\eng{kubernetes} و concurrent برابر ۱ است. برای هر Job سه container
معمول Runner یعنی build، helper و init-permissions در یک pod موقت ایجاد می‌شود.
helper image رسمی از Docker Hub استفاده شد، چون CDN مربوط به registry.gitlab
در شبکه آزمایش پاسخ ۴۰۳ می‌داد.

\subsection{Registry و Kaniko}

Registry نسخه ۲٫۸٫۳ با ClusterIP ثابت \codefile{10.96.0.50:5000} و PVC سه
گیگابایتی اجرا شد. imageهای هر commit در مسیر
\codefile{cc-phase2/<service>:<short-sha>} ذخیره می‌شوند. Kaniko جایگزین
\eng{Docker-in-Docker} شد تا Runner به socket میزبان دسترسی نداشته باشد.

image اصلی Kaniko روی gcr در این شبکه ۴۰۳ می‌داد؛ mirror
\codefile{chinayin/kaniko-project:v1.23.2-debug} استفاده شد. Registry داخلی HTTP
است، بنابراین flagهای \codefile{--insecure} و
\codefile{--insecure-registry} به executor داده شدند.

\section{طراحی GitLab CI/CD}

فایل \codefile{.gitlab-ci.yml} شامل هشت stage است:

\begin{center}
\eng{test $\rightarrow$ build $\rightarrow$ deploy $\rightarrow$ traffic
$\rightarrow$ mapreduce $\rightarrow$ validate $\rightarrow$ report
$\rightarrow$ publish}
\end{center}

\subsection{Stage تست}

در stage اول این موارد اجرا می‌شوند:

\begin{itemize}
  \item compile نحوی همه فایل‌های Python؛
  \item \codefile{bash -n} برای همه اسکریپت‌های shell؛
  \item اجرای محلی همان پنج mapper/reducer روی fixture کوچک؛
  \item تست duplicate request در Job 1؛
  \item کنترل نام متغیر و مسیر PVC سه backend؛
  \item اعتبارسنجی ۱۳ خروجی موجود.
\end{itemize}

\subsection{Stage build}

برای جلوگیری از تداخل filesystem Kaniko، هر image یک Job مستقل دارد. به دلیل
concurrent برابر یک، Jobها به ترتیب اجرا می‌شوند. team، stadium، traffic و
phase2-tools از image داخلی match همان commit به عنوان base استفاده می‌کنند تا
وابستگی به CDN خارجی کم شود.

image Hadoop در ابتدا از یک base قدیمی دو گیگابایتی ساخته می‌شد و snapshot آن
I/O زیادی ایجاد کرد. نسخه نهایی از \eng{eclipse-temurin:11-jre-jammy}،
Python و Hadoop 3.4.3 ساخته شد. دانلود Hadoop با aria2 و هشت connection انجام
شد؛ فایل ۴۹۱MB در حدود ۴ دقیقه و ۱۲ ثانیه دریافت شد. برای build این image،
memory limit Runner برابر ۲۵۰۰Mi و compressed caching غیرفعال شد.

\subsection{Stageهای deploy تا publish}

deploy manifestها را apply می‌کند و tag \codefile{dev} را با SHA commit عوض
می‌کند. traffic داده قبلی را پاک و ۱۰۰هزار درخواست تولید می‌کند. mapreduce
HDFS را scale-up و بعد از پایان scale-down می‌کند. validate همه schemaها و
عدد ۱۰۰٬۰۰۰ را کنترل می‌کند. report داشبورد را با کاربر non-root می‌سازد و
artifactهای HTML، JSON و CSV را به GitLab می‌دهد. publish Deployment وب را
restart و دریافت \codefile{summary.json} از Service داخلی را آزمایش می‌کند.

\subsection{اجرای نهایی پایپ‌لاین}

پس از اعمال آخرین اصلاح health check، commit
\codefile{f507f6c3} push شد. GitLab برای آن
پایپ‌لاین شماره ۱۶ را ساخت. صف اولیه ۱۵ ثانیه و زمان اجرای ثبت‌شده ۳۷۲ ثانیه
بود. زمان پایان در GitLab برابر ۲۵ ژوئیه ۲۰۲۶، ساعت ۰۱:۱۸ به وقت تهران ثبت شد.

\begin{table}[H]
\centering
\caption{نتیجه Jobهای پایپ‌لاین نهایی شماره ۱۶}
\small
\begin{tabularx}{\textwidth}{>{\bfseries}X p{3.0cm} p{2.4cm}}
\toprule
Job یا گروه Jobها & وضعیت & مدت تقریبی \\
\midrule
\eng{test-code} & موفق & ۶ ثانیه \\
\eng{build-match-service} & موفق & ۱۷ ثانیه \\
\eng{build-team / stadium / traffic / tools} & هر چهار موفق & ۶ تا ۸ ثانیه برای هرکدام \\
\eng{build-hadoop-streaming} & موفق & ۲۸ ثانیه \\
\eng{deploy-kubernetes} & موفق & ۱۱ ثانیه \\
\eng{generate-traffic} & موفق & ۱۱۷ ثانیه \\
\eng{run-mapreduce} & موفق & ۱۳۸ ثانیه \\
\eng{validate-output} & موفق & ۹ ثانیه \\
\eng{generate-report} & موفق & ۹ ثانیه \\
\eng{publish-report} & موفق & ۹ ثانیه \\
\bottomrule
\end{tabularx}
\end{table}

\section{Hadoop Streaming روی Kubernetes}

\subsection{راه‌اندازی HDFS با مصرف کنترل‌شده}

NameNode و DataNode به‌صورت StatefulSet تعریف شدند، اما replica پیش‌فرض آن‌ها
صفر است. تنها Job \codefile{run-mapreduce} آن‌ها را یک می‌کند. imageهای HDFS
روی میزبان موجود بودند؛ به جای دانلود دوباره دو image حدود ۲٫۰۶GB، با یک pipe
مستقیم \codefile{docker save | ctr images import} وارد containerd شدند. import
هر دو image در مجموع ۱۶٫۴ ثانیه طول کشید و فایل موقت چندگیگابایتی ساخته نشد.

پس از restart، port NameNode ممکن است Ready باشد ولی HDFS هنوز در Safe Mode
باشد. در یکی از اجراها upload با پیام
\eng{Cannot create file; Name node is in safe mode} متوقف شد. اسکریپت نهایی
تا ۱۸۰ ثانیه، هر دو ثانیه وضعیت \codefile{hdfs dfsadmin -safemode get} را
می‌خواند و فقط پس از OFF شدن نوشتن را شروع می‌کند.

\subsection{Job 1: Parsing and Cleaning}

چهار فایل خام به HDFS فرستاده می‌شوند. mapper JSON را parse می‌کند، فیلدهای
اجباری و بازه status code و زمان را می‌سنجد و رکورد را به nginx، service یا
invalid tag می‌دهد. reducer با کلید hash شده request ID، retryهای تکراری را
حذف می‌کند. خروجی‌های این مرحله دو CSV تمیز و فایل invalid هستند.

\subsection{Job 2: General Nginx Aggregation}

برای service، endpoint و scenario تعداد کل، موفق، خطای ۴xx، خطای ۵xx، نرخ خطا
و میانگین response time محاسبه می‌شود. همین خروجی منبع summary نهایی است.

\subsection{Job 3: Country--Entity Count}

service log تمیز به سه دسته team، match day و stadium/city نگاشت می‌شود و تعداد
درخواست هر موجودیت برای هر کشور جمع می‌شود.

\subsection{Job 4: Popular Entity by Country}

خروجی Job 3 دوباره پردازش می‌شود. برای هر کشور محبوب‌ترین تیم، روز بازی و
ورزشگاه انتخاب می‌شود. در تساوی، مرتب‌سازی نام نتیجه را deterministic می‌کند.

\subsection{Job 5: Final Summary}

Job آخر خروجی‌های Jobهای ۲ تا ۴ را با metadata پیش‌بینی ترکیب می‌کند و
\codefile{outputs/final/summary.json} را می‌سازد. پس از آن script validator
ساختار همه خروجی‌ها را کنترل می‌کند.

\begin{resultbox}{اجرای واقعی پنج Job}
در پایپ‌لاین نهایی، Job MapReduce از شروع scale-up تا پایان تحلیل ۱۳۸ ثانیه
طول کشید. pod پردازش هیچ restart یا OOM نداشت. هر پنج marker اجرایی، اعتبارسنجی
۱۳ فایل و پیام \eng{MapReduce pipeline completed successfully} در trace ثبت شد.
\end{resultbox}

\section{خروجی و تحلیل نتایج}

\begin{table}[H]
\centering
\caption{آمار سرویس‌ها در اجرای ۱۰۰هزار درخواست}
\small
\begin{tabular}{lrrrrr}
\toprule
\eng{service} & کل & موفق & ۴xx & ۵xx & میانگین ms \\
\midrule
\eng{match-service} & ۲۸٬۱۶۲ & ۲۷٬۲۹۲ & ۵۶۳ & ۳۰۷ & ۳۲٫۳۱۱ \\
\eng{stadium-service} & ۲۲٬۲۱۹ & ۲۱٬۵۶۵ & ۴۵۱ & ۲۰۳ & ۳۲٫۸۳۵ \\
\eng{team-service} & ۴۹٬۶۱۹ & ۴۸٬۱۵۸ & ۹۷۰ & ۴۹۱ & ۲۶٫۰۵۳ \\
\bottomrule
\end{tabular}
\end{table}

team-service با ۴۹٬۶۱۹ درخواست پرتقاضاترین سرویس بود. نرخ خطای match-service
برابر ۳٫۰۸۹۳ درصد و از دو سرویس دیگر کمی بیشتر بود. کندترین endpoint طبق
summary، \codefile{/api/stadiums} شد.

\begin{table}[H]
\centering
\caption{آمار scenarioها}
\small
\begin{tabular}{lrrrr}
\toprule
\eng{scenario} & تعداد & موفق & نرخ خطا & میانگین ms \\
\midrule
\eng{normal} & ۹۱٬۴۶۶ & ۹۱٬۴۶۶ & ۰\% & ۲۵٫۶۸۴ \\
\eng{popular} & ۵٬۰۷۲ & ۵٬۰۷۲ & ۰\% & ۲۵٫۳۰۷ \\
\eng{invalid} & ۱٬۹۸۴ & ۰ & ۱۰۰\% & ۲۵٫۵۰۷ \\
\eng{server\_error} & ۱٬۰۰۱ & ۰ & ۱۰۰\% & ۲۵٫۶۵۲ \\
\eng{slow} & ۴۷۷ & ۴۷۷ & ۰\% & ۷۹۳٫۳۰۰ \\
\bottomrule
\end{tabular}
\end{table}

سناریوی slow فقط ۰٫۴۷۷ درصد ترافیک بود، ولی میانگین زمان آن حدود ۷۹۳ms شد؛
پس میانگین کلی endpointها را بدون ایجاد load کنترل‌نشده افزایش داد.

summary نهایی این نتایج را ثبت کرد:

\begin{itemize}
  \item محبوب‌ترین تیم کلی: آرژانتین؛
  \item پرتقاضاترین روز بازی: ۲۵ ژوئن ۲۰۲۶؛
  \item پرتقاضاترین ورزشگاه: New York New Jersey Stadium؛
  \item قهرمان پیش‌بینی‌شده: آرژانتین؛
  \item فینال پیش‌بینی‌شده: فرانسه مقابل آرژانتین؛
  \item محل فینال پیش‌بینی‌شده: New York New Jersey Stadium.
\end{itemize}

\section{Dashboard و انتشار نتیجه}

اسکریپت \codefile{report-web/generate\_dashboard.py} summary و service stats را
می‌خواند و HTML مستقل تولید می‌کند. image ابزار با UID برابر ۱۰۰۰۱ اجرا می‌شود.
در اولین اجرای report، پوشه \codefile{outputs/final} توسط Hadoop با مالک root
و mode برابر ۰۷۵۵ ساخته شده بود و نوشتن index با PermissionError متوقف شد.

نسخه نهایی یک initContainer بسیار کوچک دارد که فقط group همان پوشه را ۲۰۰۰ و
mode را ۰۷۷۵ می‌کند. container اصلی همچنان non-root باقی می‌ماند. smoke test
موفق پیام زیر را ثبت کرد:

\begin{latin}
\begin{lstlisting}
Dashboard written to /project-data/outputs/final/index.html
<title>World Cup Log Analytics
\end{lstlisting}
\end{latin}

Report Web فایل‌های HTML و JSON را read-only mount می‌کند. دسترسی میزبان از
\codefile{http://localhost:9000/index.html} و دسترسی داخلی از Service
\codefile{report-web} است.

\section{کنترل RAM، CPU و I/O}

\begin{table}[H]
\centering
\caption{اقدام‌های انجام‌شده برای کنترل منابع}
\small
\begin{tabularx}{\textwidth}{>{\bfseries}p{4.1cm}X}
\toprule
بخش & اقدام \\
\midrule
Runner & concurrent برابر ۱؛ buildهای مستقل و ترتیبی \\
GitLab & خاموش‌کردن monitoring غیرضروری و تعیین request/limit \\
Kaniko & cache Registry، base داخلی و compressed caching خاموش فقط برای Hadoop \\
Traffic & ۳۲ worker محدود و سقف ۲ CPU و ۱۹۲MiB \\
HDFS & replica صفر در حالت عادی و scale موقت در MapReduce \\
Validator & خواندن CSV به‌صورت streaming به جای list کامل \\
Image انتقالی & pipe مستقیم به containerd و عدم ساخت archive موقت \\
Spark & توقف GitLab و backendها هنگام اجرای بخش اختیاری \\
\bottomrule
\end{tabularx}
\end{table}

کمترین حافظه available مشاهده‌شده در اجرای Hadoop حدود ۱٫۰ تا ۱٫۲GiB بود.
هیچ pod پردازش restart نشد. نسخه اول validator با limit برابر ۱۲۸MiB
OOMKilled شد، چون دو CSV بزرگ را کامل در list می‌ریخت. نسخه streaming روی
خروجی ۱۰۰هزار رکورد با tracemalloc حدود ۰٫۹MiB peak تخصیص Python داشت و همان
۱۳ بررسی schema را حفظ کرد.

\section{دفترچه خطاها و اصلاح مرحله‌به‌مرحله}

خطاهای زیر در اجرای واقعی دیده شدند. این موارد حذف یا پنهان نشده‌اند، چون بخش
مهم فاز دوم خطایابی محیط واقعی Kubernetes و CI/CD است.

\begin{enumerate}
  \item \textbf{Load مستقیم kind:} image چندسکویی در Docker 29 و containerd 2
  digest ناقص داد. برای imageهای معمول pull مستقیم و برای imageهای HDFS import
  مستقیم با ctr استفاده شد.

  \item \textbf{Readiness GitLab:} آدرس loopback و مسیر curl اولیه درست نبود.
  probe به curl داخلی GitLab و endpoint رسمی readiness تغییر کرد.

  \item \textbf{Runner helper:} دانلود از CDN مربوط به registry.gitlab پاسخ
  ۴۰۳ داشت. helper رسمی همان نسخه از Docker Hub انتخاب شد.

  \item \textbf{کلیدهای GitLab:} config موقت باعث تغییر کلید encryption پس از
  تعویض pod شد. PVC جداگانه gitlab-config اضافه شد.

  \item \textbf{Image ابزار CI:} tag انتخابی bitnami/kubectl وجود نداشت.
  image \codefile{alpine/k8s:1.32.3} جایگزین شد.

  \item \textbf{Kaniko روی gcr:} دسترسی ۴۰۳ بود؛ mirror Docker Hub استفاده شد.

  \item \textbf{Build اولیه Hadoop:} base قدیمی دو گیگابایتی I/O بسیار زیاد
  ایجاد کرد و یک بار EOF و یک بار OOM رخ داد. base سبک Java 11 و Hadoop 3.4.3
  ساخته شد و limit همان Job به ۲۵۰۰MiB رسید.

  \item \textbf{دانلود Hadoop:} curl تک‌connection سرعت حدود ۰٫۳MB/s داشت.
  aria2 با هشت connection میانگین حدود ۱٫۹MiB/s داد.

  \item \textbf{RBAC namespace:} Runner مجوز get namespace نداشت. ClusterRole
  فقط برای نام cc-project افزوده شد و list همچنان غیرمجاز ماند.

  \item \textbf{تداخل buildها:} اجرای شش Kaniko پشت سر هم در یک pod باعث شد
  بعضی imageها فایل‌های base را درست نداشته باشند. هر image به Job مستقل تبدیل
  شد.

  \item \textbf{دسترسی Registry HTTP:} pull داخلی ابتدا HTTPS فرض شد.
  flag insecure و تنظیم mirror اضافه شد.

  \item \textbf{DNS Registry در node:} containerd به DNS pod دسترسی نداشت.
  IP ثابت Service در hosts node و اسکریپت ایجاد کلاستر ثبت شد.

  \item \textbf{مسیر access log:} Nginx در شاخه‌ای متفاوت از ورودی Hadoop
  می‌نوشت. MapReduce متوقف، مسیر اصلاح، فایل اشتباه حذف و ترافیک از نو تولید شد.

  \item \textbf{چهار retry شبکه‌ای:} در یک اجرای آزمایشی ۱۰۰٬۰۰۴ خط خام و
  چهار request ID تکراری دیده شد. dedup مبتنی بر request ID در Job 1 و تست
  رگرسیون اضافه شد.

  \item \textbf{متغیر لاگ backend:} LOG\_PATH توسط برنامه خوانده نمی‌شد.
  SERVICE\_LOG\_PATH و مسیر کامل PVC برای هر سه سرویس اصلاح شد.

  \item \textbf{Validator OOM:} list کردن CSVهای بزرگ بیش از ۱۲۸MiB مصرف کرد.
  validator سطر‌به‌سطر بازنویسی شد.

  \item \textbf{HDFS Safe Mode:} readiness شبکه پیش از امکان نوشتن برقرار شد.
  انتظار صریح حداکثر ۱۸۰ ثانیه برای Safe Mode OFF اضافه شد.

  \item \textbf{ثبت IP جدید DataNode:} پس از ساخت دوباره pod، NameNode به‌علت
  نبود reverse-DNS، IP جدید DataNode را رد کرد. hostname check مخصوص محیط
  container غیرفعال شد؛ احراز هویت و شبکه namespace تغییری نکرد.

  \item \textbf{Permission داشبورد:} Hadoop پوشه final را root:root ساخته بود.
  initContainer فقط group و write permission همان پوشه را آماده کرد.

  \item \textbf{اتصال هنگام rollout گزارش:} اولین wget به endpoint در حال حذف
  گیر کرد، در حالی که درخواست بعدی پاسخ ۲۰۰ داشت. برای health check انتشار،
  timeout پنج‌ثانیه‌ای و retry محدود به فرمان wget اضافه شد.

  \item \textbf{خاموش‌شدن لپ‌تاپ میان اجرا:} اجرای شماره ۱۴ هنگام تولید ترافیک
  با قطع برق میزبان نیمه‌کاره ماند. پس از روشن‌شدن دوباره، کلاستر kind،
  PVCهای GitLab، Registry و داده پروژه بررسی شدند و سالم بودند. Job نیمه‌کاره
  مبنای نتیجه قرار نگرفت؛ pipeline تازه از آخرین commit اجرا شد و اجرای شماره
  ۱۶ از test تا publish کامل و موفق پایان یافت.

  \item \textbf{تفکیک شواهد دو فاز:} در بازبینی نهایی پیش از انتشار GitHub
  مشخص شد پوشه evidence فاز دوم هنوز تصویرها و متن‌های اجرای Docker Compose
  فاز اول را نگه داشته است. آن فایل‌ها برای اثبات Kubernetes مناسب نبودند و
  از فاز دوم کنار گذاشته شدند؛ نسخه اصلی آن‌ها در گزارش فاز اول محفوظ ماند.
  پوشه evidence فاز دوم با وضعیت واقعی namespace، اطلاعات Pipeline 16،
  شمارش لاگ‌ها، خروجی MapReduce، publish و Spark روی Kubernetes از نو ساخته شد.
\end{enumerate}

\section{بخش امتیازی Spark روی Kubernetes}

\begin{scorebox}{طراحی بخش اختیاری}
کد Structured Streaming، manifest اجرای Spark و Job خروجی‌دادن batchها در
پوشه‌های \codefile{spark/} و \codefile{k8s/spark/} قرار دارند. ورودی و خروجی
هر دو از PVC مشترک استفاده می‌کنند و اجرای Spark مستقل از مسیر اجباری
GitLab/Hadoop است.
\end{scorebox}

اسکریپت \codefile{run\_spark\_kubernetes.sh} برای کنترل منابع ابتدا GitLab،
Runner، Nginx و backendها را scale-down می‌کند. ConfigMap برنامه Spark را
می‌سازد، Job Spark را اجرا می‌کند، batchهای اتمیک را از لاگ‌های
واقعی صادر می‌کند و پس از complete شدن هر دو Job، log نتیجه را نمایش می‌دهد.

\begin{resultbox}{نتیجه اجرای واقعی Spark}
image رسمی \eng{apache/spark:3.5.5} بدون فایل موقت و با pipe مستقیم وارد
containerd شد. Job با \eng{local[2]}، حافظه driver برابر ۷۶۸MiB و limit کل
۱۶۰۰MiB اجرا شد و در ۶۸ ثانیه، بدون restart، کامل شد. صادرکننده پنج batch
دویست‌سطری Nginx و پنج batch دویست‌سطری service ساخت. Structured Streaming
در پنجره ۱۰ثانیه‌ای ۴۹۴ درخواست team، ۲۸۰ درخواست match و ۲۲۶ درخواست stadium
را در اولین batch دید و تیم محبوب را برای هشت کشور محاسبه کرد. برای نمونه،
آرژانتین با ۱۱۳ مشاهده برای کشور آرژانتین و ژاپن با ۸۱ مشاهده برای کشور ژاپن
در خروجی زنده ثبت شدند. پس از پایان، trap اسکریپت GitLab، Runner، Nginx و هر
سه backend را بازیابی کرد و rollout همه آن‌ها موفق بود.
\end{resultbox}

\section{آزمون‌ها و معیار پذیرش}

\begin{table}[H]
\centering
\caption{ماتریس آزمون فاز دوم}
\small
\begin{tabularx}{\textwidth}{X p{2.4cm} X}
\toprule
آزمون & نتیجه & شاهد \\
\midrule
سلامت node و namespace & موفق & node برابر Ready و namespace برابر cc-project \\
GitLab داخل Kubernetes & موفق & Deployment Ready و صفحه ورود روی ۸۹۲۹ \\
Runner Kubernetes executor & موفق & pod موقت برای هر Job و Runner online \\
build با Kaniko & موفق & شش image با tag commit در Registry خصوصی \\
deploy سرویس‌ها & موفق & rollout چهار Deployment اصلی \\
ترافیک فقط از Nginx & موفق & مقصد Job فقط nginx-gateway و ۱۰۰هزار ID یکتا \\
لاگ مشترک & موفق & ۱۰۰هزار Nginx و ۱۰۰هزار service log روی PVC \\
پنج Job Hadoop & موفق & پنج marker و پیام completion \\
اعتبارسنجی & موفق & ۱۳ فایل و total برابر ۱۰۰٬۰۰۰ \\
Dashboard & موفق & HTML، JSON و NodePort 9000 \\
محدودیت منابع & موفق & صفر restart در Hadoop و validator با حافظه ثابت \\
Spark Structured Streaming & موفق & ده batch، پنجره زنده، زمان ۶۸ ثانیه و صفر restart \\
\bottomrule
\end{tabularx}
\end{table}

\section{شواهد نهایی قابل بررسی بدون دسترسی به لپ‌تاپ}

GitLab این پروژه داخل کلاستر محلی اجرا شده است و پس از تحویل، مصحح نمی‌تواند
آدرس \codefile{localhost:8929} لپ‌تاپ را باز کند. به همین دلیل فقط به تصویر
صفحه GitLab اکتفا نشد. وضعیت نهایی با API خود GitLab و فرمان‌های Kubernetes
خوانده، از Secretها پاک و در پوشه \codefile{report/evidence} ذخیره شد. تمام
فایل‌ها متنی هستند تا در GitHub مستقیماً دیده و مقایسه شوند.

\begin{table}[H]
\centering
\caption{فهرست شواهد واقعی فاز دوم}
\small
\begin{tabularx}{\textwidth}{>{\bfseries}p{4.7cm}X}
\toprule
فایل & محتوای قابل بررسی \\
\midrule
\codefile{01\_kubernetes\_status.txt} & Deploymentها، StatefulSetها، Serviceها و PVCهای namespace \\
\codefile{02\_gitlab\_pipeline\_16.txt} & commit، زمان کل و وضعیت و مدت هر ۱۳ Job \\
\codefile{03\_traffic\_and\_outputs.txt} & شمارش لاگ، شناسه یکتا و summary نهایی \\
\codefile{04\_mapreduce\_and\_validation.txt} & هفت گام اجرا، پنج Job تحلیلی و اعتبارسنجی ۱۳ فایل \\
\codefile{05\_dashboard\_and\_publish.txt} & rollout وب، پاسخ‌های HTTP 200 و نتیجه publish \\
\codefile{06\_spark\_kubernetes.txt} & وضعیت Job، پنجره زنده، limit حافظه و بازیابی سرویس‌ها \\
\codefile{07\_security\_and\_reproduction.txt} & کنترل Secret و مسیر بازتولید کامل سامانه \\
\bottomrule
\end{tabularx}
\end{table}

\begin{latin}
\begin{lstlisting}
Pipeline 16 | commit f507f6c3 | status success | duration 372 s
traffic-generator:
  Completed: 100,000 in 98.30s (1017.3 req/s)
data checks:
  nginx lines=100000, unique request IDs=100000
  backend lines=28162+49619+22219=100000
mapreduce:
  five analysis jobs completed
  Validated 13 required output files.
publish:
  dashboard HTTP=200, summary HTTP=200
spark:
  status=Complete, duration=68s, restarts=0
\end{lstlisting}
\end{latin}

\begin{resultbox}{قابلیت ممیزی مخزن GitHub}
کد، manifestها، اسکریپت نصب، داده واقعی، خروجی‌های تحلیل، گزارش PDF و شواهد
اجرای نهایی در یک commit قرار گرفته‌اند. بزرگ‌ترین فایل ۲۸٫۶MiB است، بنابراین
هیچ فایل از محدودیت انتشار معمول GitHub عبور نمی‌کند. مصحح می‌تواند بدون
دسترسی به Secretها نتیجه را بررسی کند و در صورت نیاز کل محیط را از روی README
بازتولید کند.
\end{resultbox}

\section{روش اجرای پروژه}

\subsection{پیش‌نیاز}

ابزارهای Docker، kind، kubectl، git، curl و openssl لازم هستند. کلاستر با این
فرمان ساخته می‌شود:

\begin{latin}
\begin{lstlisting}[language=bash]
bash scripts/phase2/create_cluster.sh
bash scripts/phase2/load_core_images.sh
\end{lstlisting}
\end{latin}

\subsection{راه‌اندازی GitLab و Runner}

\begin{latin}
\begin{lstlisting}[language=bash]
bash scripts/phase2/deploy_gitlab.sh
bash scripts/phase2/configure_gitlab.sh
bash scripts/phase2/wait_for_pipeline.sh
\end{lstlisting}
\end{latin}

اسکریپت configure token کوتاه‌عمر می‌سازد، پروژه را ایجاد می‌کند، Runner را
ثبت می‌کند و مخزن را push می‌کند. token در خروجی چاپ یا commit نمی‌شود.

\subsection{کنترل وضعیت و آدرس‌ها}

\begin{latin}
\begin{lstlisting}[language=bash]
bash scripts/phase2/collect_status.sh
kubectl -n cc-project get pods,svc,jobs,pvc
\end{lstlisting}
\end{latin}

\begin{itemize}
  \item GitLab: \codefile{http://localhost:8929}
  \item Nginx API: \codefile{http://localhost:8080}
  \item Dashboard: \codefile{http://localhost:9000/index.html}
\end{itemize}

\subsection{اجرای Spark اختیاری}

\begin{latin}
\begin{lstlisting}[language=bash]
IMAGE_TAG=<successful-short-sha> \
  bash scripts/phase2/run_spark_kubernetes.sh
\end{lstlisting}
\end{latin}

\section{ساختار فایل‌های تحویلی}

\begin{latin}
\begin{lstlisting}
code/
  .gitlab-ci.yml
  k8s/
    base/        # namespace, PVC, Registry, backend, Nginx, HDFS, report
    jobs/        # traffic, reset, MapReduce, validator, report
    gitlab/      # GitLab CE
    runner/      # RBAC and Runner
    spark/       # optional Spark workload
  match-service/
  team-service/
  stadium-service/
  traffic-generator/
  mapreduce/     # five mapper/reducer pairs
  hadoop-k8s/
  phase2-tools/
  report-web/
  scripts/phase2/
  tests/
  docs/architecture.{dot,svg,png}
report/
  report.tex
  report.pdf
  architecture.png
  evidence/       # sanitized Kubernetes, GitLab CI, Hadoop and Spark proof
\end{lstlisting}
\end{latin}

نام بسته نهایی مطابق دستورکار به‌صورت زیر است:

\begin{center}
\codefile{CC\_Project\_Phase2\_40231901\_40231064.zip}
\end{center}

\section{جمع‌بندی}

در این فاز فقط manifest نوشته نشد؛ GitLab و Runner واقعی داخل همان کلاستر بالا
آمدند، مخزن واقعی push شد و pipeline از build تا publish اجرا شد. مسیر داده
از یک Job ترافیک ۱۰۰هزار درخواستی، Nginx و سه backend شروع شد، با HDFS و پنج
Job Streaming ادامه پیدا کرد و به ۱۳ خروجی تأییدشده و dashboard قابل مشاهده
رسید.

خطاهای اصلی پروژه بیشتر از جنس اتصال اجزا بودند: DNS بین host namespace و
cluster، کلیدهای پایدار GitLab، مسیر دقیق PVC، نام متغیر environment، Safe Mode
HDFS، مصرف حافظه validator و permission کاربر non-root. هر مورد پس از مشاهده
در اجرای واقعی به کد یا تست رگرسیون تبدیل شد. در نتیجه نسخه تحویلی صرفاً یک
نمونه نمایشی نیست و مسیر نصب و اجرای آن از روی اسکریپت‌ها قابل تکرار است.

\end{document}
